\begin{table}[H]
\centering
\begin{threeparttable}
\caption{Illustrative Construction of an English-Learning Module from a Source Article}
\label{tab:module_construction}
\scriptsize
\setlength{\tabcolsep}{4pt}
\begin{tabularx}{\textwidth}{p{1.9cm}p{4.0cm}X}
\toprule
Stage & Standard transformation rule & Worked example using CTRL\_021 (\textit{The Economist}, ``The lessons from the brazen heist at the Louvre'') \\
\midrule
Source article & Select a blocked foreign-source article that fits the bank definition, has enough factual density for close reading, and passes basic classroom-safety and platform-compliance review. Items judged likely to contain graphic violence, explicit self-harm, illegal instruction, or other content that could violate platform rules or create reader-safety risks are excluded rather than adapted. & The article is a non-China culture report on the October 2025 Louvre jewelry theft. It sits in the non-China control pool because it contains no China-related political content, while still offering a clear narrative, institutional stakes, and advanced vocabulary suitable for TOEFL / IELTS / GRE / GMAT preparation. \\
\addlinespace
Abridged reading passage & Remove site navigation, captions, and promotional text; preserve chronology, main actors, causal logic, and one analytic takeaway; restructure the piece into a passage that can be read and annotated inside the app. & The source report is converted into a roughly 1,000-word classroom passage organized around five moves: the raid, why the Louvre matters, how museum thefts usually work, what failed in this case, and the broader security lesson that museums must slow thieves down rather than rely on alarms alone. The app version shortens the original, simplifies transitions, and keeps the classroom focus on comprehension and reusable language rather than full article reproduction. \\
\addlinespace
Vocabulary and syntax notes & Highlight 5--8 collocations or discipline-specific terms and pair them with one or two sentence-level structures for explanation, translation, and reuse. & Example vocabulary targets include \textit{make off with}, \textit{strike at the heart of}, \textit{dismantle the pieces}, \textit{pose serious questions about}, and \textit{buy time}. Example syntax targets include a \textit{that}-clause subject (``That a truck could park outside the museum ...'') and the contrastive frame \textit{not only ... but also ...}, matching the phrase-based annotation format shown in the app screenshots. \\
\addlinespace
Video / audio explainer & Add a short instructor-led recap that reinforces the article's factual sequence, clarifies one inferential takeaway, and previews the language focus for the lesson. & A 60--90 second explainer for this module walks through the seven-minute theft, explains why museum-security specialists focus on delay rather than perfect prevention, and pauses over highlighted collocations and sentence patterns on the teaching slide before the quiz stage. \\
\addlinespace
Quiz battery & Convert the passage into multiple item types that test comprehension and language uptake rather than political agreement. The teaching-data structure scores five subtypes from 0 to 20 and sums them to a 0--100 total. & The quiz follows the app's observed categories: reading comprehension (yuedu lijie), grammar / syntax (yufa jufa), vocabulary (danci ceshi), listening (tingli ceshi), and sentence accumulation (juzi jilei). In the standardized scoring template, each subtype contributes 20 points, so the full lesson quiz totals 100. Illustrative items include: why experts treat the Louvre case as unusually serious; what \textit{make off with} means in context; what the \textit{that}-clause contributes in a sentence about security failure; what the instructor identifies as the main policy lesson; and how to reconstruct a model sentence such as ``The goal of security is to delay thieves long enough for a response.'' \\
\addlinespace
Weekly evaluation & End each module with a short standardized evaluation of engagement and future demand. & After completing the CTRL\_021 lesson, participants would rate how interesting the lesson was, how credible it seemed, and whether they wanted similar materials in future weeks. In the experiment, these weekly evaluations form the bridge between attention to a lesson and later demand for related content. \\
\bottomrule
\end{tabularx}
\begin{tablenotes}[flushleft]
\footnotesize
\item Note: This table documents the pedagogical workflow rather than reproducing copyrighted source text. The worked example uses one article from the finalized non-China control bank to show how a source report would be screened, adapted, and converted into the study's standardized module format.
\end{tablenotes}
\end{threeparttable}
\end{table}
